\emph{The mechanics of radical letters.} Work in the blade basis (Fact \claimref{conformal.blade-basis-persists-under-degeneracy}) adapted to a decomposition $W = V_1 \oplus R$: each blade is $x_T\, \rho_S$ with $x_T$ a blade in a basis of $V_1$ and $\rho_S := \rho_{s_1} \cdots \rho_{s_j}$, $S = \{s_1 < \dots < s_j\}$. A radical letter $\rho$ satisfies $B(\rho, z) = 0$ for every vector $z$, so $\rho z = -z\rho$ with no scalar term: in any product of blades, radical letters reorder by pure anticommutation and contract with nothing, while a repeated radical letter makes the product zero ($\rho^2 = 0$). Hence
\[
(x_T\, \rho_S)(x_{T'}\, \rho_{S'}) \;=\;
\begin{cases}
0 & \text{if } S \cap S' \neq \emptyset, \\
\pm\, (x_T\, x_{T'})\; \rho_{S \cup S'} & \text{if } S \cap S' = \emptyset ,
\end{cases}
\]
where $x_T\, x_{T'}$ expands over blades of $V_1$ alone: radical letters accumulate and are never contracted away.

The ideal, its dimension and its nilpotency index are parts (i) and (ii) of Theorem \claimref{std-algebra.radical-ideal-and-quotient} for the space of signature $(p, q, r)$; in the blade notation above they read as follows.

\emph{The ideal, and its dimension.} $\mathcal{S} := \operatorname{span}\{x_T \rho_S : S \neq \emptyset\}$ is a two-sided ideal by the display, and it equals $J$: every $x_T \rho_S$ with $s_1 \in S$ lies in $\Cl \cdot \rho_{s_1} \cdot \Cl$, and conversely every product $a \rho_i b$ expands into blades each retaining the letter $\rho_i$. Counting blades with $S \neq \emptyset$ gives $\dim J = 2^m - 2^{m-r}$.

\emph{Nilpotency, and its exact index.} By the display, a product of $j$ elements of $J$ lies in the span of blades carrying at least $j$ radical letters; only $r$ radical letters exist, so $J^{\,r+1} = 0$. In the other direction, $\rho_1 \cdots \rho_r$ is a product of $r$ elements of $J$ and is a basis blade, hence nonzero: $J^{\,r} \neq 0$. For $r \geq 2$ the nonzero blade $\rho_1 \rho_2 \in J^2$ shows the ideal is not square-zero.

\emph{The unit ideal in the completion.} Let $(\widetilde W, Q)$ be any completion and $\rho_i$ a radical basis vector. By nondegeneracy there is $\varphi_i \in \widetilde W$ with $Q(\rho_i, \varphi_i) = 1$, so the Clifford relation (\claimref{std-algebra.vector-squares}) gives $\rho_i \varphi_i + \varphi_i \rho_i = 2$, and $1 = \tfrac12(\rho_i \varphi_i + \varphi_i \rho_i)$ lies in the two-sided ideal generated by $\rho_i$ (both summands carry the factor $\rho_i$). Replacing $\varphi_i$ by $\varphi_i - \tfrac12 Q(\varphi_i, \varphi_i)\, \rho_i$ leaves $Q(\rho_i, \varphi_i) = 1$ (since $Q(\rho_i, \rho_i) = 0$) and makes $\varphi_i$ null, so $\operatorname{span}\{\rho_i, \varphi_i\}$ is a hyperbolic plane. This is the finite-dimensional instance of the resolution mechanism of \cite[Prop.~3.4]{ButlerT4}.
